\begin{figure}[H]
\centering
\includegraphics[width=1.0\linewidth]{fidelity/G7-01-wakayama_trajectory_map.pdf}
\caption[冬の 1 例]{2020 年 1 月 15 日 00 UTC に和歌山の中心格子から上げた場合の飛行経路。左が全体、右が着地点の拡大。地点のみ固定（橙の破線）と両方固定（黄）が南側にそろい、基準（黒）と時刻のみ固定（青の破線）が北側にそろう。冬は地点の固定による誤差のほうが大きいことが、1 例からも見てとれる。}
\label{fig:g7-01-wakayama}
\end{figure}

\begin{figure}[H]
\centering
\includegraphics[width=1.0\linewidth]{fidelity/G7-07-wakayama_trajectory_map.pdf}
\caption[夏の 1 例]{同じ図の 7 月版。飛行距離が \SI{80}{\kilo\metre} ほどと短く、4 通りの差も \SI{2.5}{\kilo\metre} にとどまる。拡大パネルの縮尺が冬（\SI{14}{\kilo\metre}）と夏（\SI{2}{\kilo\metre}）で違う点に注意。}
\label{fig:g7-07-wakayama}
\end{figure}
